\section{Conclusions and Future Work}
\label{sec:conclusions}

\subsection{Conclusions}

This study investigated whether information extracted by a
Self-Organizing Map (SOM) could provide useful guidance for adaptive
mesh refinement in CFD. The two-dimensional lid-driven cavity at
\(Re=100\) was used as a controlled test problem, and all proposed
refinement strategies were compared at identical refinement budgets.

Before introducing machine learning, a uniform-mesh convergence study
was performed using M20, M40, M80, and M160 meshes. The successive
centerline differences decreased substantially with refinement, with
observed convergence rates approaching second-order behavior on the
finer mesh levels. The M160 solution was subsequently used as a
fine-grid numerical reference rather than as an exact solution.

A \(4\times4\) SOM was trained using flow information extracted from the
M40 pilot solution. Spatial coordinates were deliberately excluded from
the SOM input. Nevertheless, the trained map produced spatially coherent
flow-state regions when its classifications were projected back onto the
physical cavity. This demonstrates that the SOM learned structured
relationships among the local flow variables rather than simply
encoding spatial position.

SOM quantization error initially appeared potentially useful because it
was concentrated in regions that also exhibited relatively large
M40--M160 discrepancies. However, the controlled adaptive CFD
experiments showed that this correspondence was insufficient to make
quantization error an effective direct refinement indicator.

At all four tested refinement budgets, the conventional velocity-gradient
criterion produced a lower RMS discrepancy than direct SOM-QE refinement.
Gradient refinement also improved monotonically with increasing budget,
whereas SOM-QE exhibited non-monotonic behavior between the 5\% and 10\%
budgets.

The final adaptive meshes demonstrate that this difference cannot be
attributed simply to the number of cells. At the 10\% budget, both
SOM-QE and Gradient produce meshes containing 2080 cells, but they
allocate those cells differently. The Gradient strategy concentrates
more resolution near regions of strong velocity variation, while SOM-QE
allocates part of the same finite refinement budget to additional flow
states that are unusual in SOM feature space.

The results therefore establish an important distinction:

\begin{equation}
\text{representation novelty}
\neq
\text{discretization error}
\neq
\text{refinement utility}.
\end{equation}

A cell can be unusual from the perspective of an unsupervised
machine-learning model without being the location where additional
spatial resolution provides the greatest numerical benefit.

The refinement-benefit analysis reinforces this conclusion. Refining
cells according to the Gradient criterion improved the solution over
regions extending beyond the directly selected cells, illustrating the
global coupling of the CFD solution. Consequently, local indicator
quality cannot be assessed solely through correlation with a pre-existing
local discrepancy field.

A second SOM-derived quantity, the regime-transition indicator \(T\),
was subsequently investigated. Unlike quantization error, this quantity
was only weakly correlated with the velocity-gradient magnitude and
therefore represented more distinct information.

This motivated the Hybrid-v1 indicator,

\begin{equation}
I_{\mathrm{hyb}}
=
G^{*}(1+0.5T),
\end{equation}

in which the conventional Gradient criterion remained dominant and SOM
information only modified the ranking of cells near the refinement
threshold.

Hybrid-v1 behaved as intended. Approximately 91--96\% of the
Gradient-selected cells were preserved, while a small number of
high-transition cells replaced cells having somewhat larger gradients
but little SOM-transition activity.

At the 10\% budget, this modification reduced the RMS discrepancy by
approximately 1.24\% relative to Gradient. However, the improvement was
not reproduced at the 5\%, 15\%, or 20\% budgets.

Hybrid-v1 therefore does not demonstrate a systematic advantage over
Gradient refinement.

The principal result of this investigation is consequently not that
machine learning improves CFD mesh refinement. Instead, the study
demonstrates that SOM-derived flow-state information can be physically
structured and complementary to a conventional numerical indicator,
while also showing that converting such information into reliable
refinement decisions is a substantially more difficult problem.

For the present \(Re=100\) cavity experiments, the velocity-gradient
criterion remains the most consistently accurate refinement strategy.


\subsection{Second-stage SOM-interface result}

The negative SOM-QE result led to a reformulation of the machine-learning
hypothesis. Rather than treating quantization error as a direct error
indicator, the SOM BMU field was used to identify interfaces between
different learned flow states. A new ranking,

\begin{equation}
I_{\mathrm{SI}}
=
T\sqrt{G^*\Delta U^*},
\end{equation}

combined the SOM transition fraction with normalized velocity-gradient
magnitude and the local velocity-vector jump across the interface.

The criterion was evaluated using 40 fixed interface seeds and one- or
two-layer graph expansion. The resulting selections contained 78 and 114
cells. Conventional Gradient controls were then constructed with exactly
the same selected-cell counts, producing final meshes of 1834 and 1942
cells.

In these cell-matched comparisons, SOM-interface reduced the volume-weighted
mean velocity discrepancy by approximately \(13.9\%\) and \(18.7\%\), and
the RMS discrepancy by approximately \(1.4\%\) and \(4.0\%\), for the \(1N\)
and \(2N\) cases, respectively.

These results represent a change in the evidence compared with the earlier
SOM-QE experiments. They do not show that SOM information is intrinsically
superior to gradient-based refinement. Instead, they show that the
\emph{type} of SOM information used matters: direct representation error
performed poorly, whereas spatial transitions between learned flow states,
combined with local physical variation, produced a different and promising
refinement pattern.

The cell-selection overlap was \(67.95\%\) for \(1N\) and \(64.91\%\) for
\(2N\). Thus, the observed differences arose from a substantial but not
complete change in the spatial allocation of the refinement budget.

The refinement-utility analysis further showed that SOM-interface does not
uniformly outperform Gradient in every spatial subset. Gradient retained
lower mean post-refinement error in the common and Gradient-only regions,
while SOM-interface showed lower error in the SOM-only region for the
\(2N\) case and in the large region selected by neither method. The global
improvement is therefore interpreted as a property of the coupled CFD
solution rather than as evidence that every SOM-selected cell has higher
individual utility.

The appropriate conclusion at this stage is that the SOM-interface
hypothesis has passed an initial controlled test and merits further
investigation, but requires validation across additional flow conditions,
geometries, and independent experiments before any general claim can be
made.


\subsection{SOM architecture sensitivity}

A complementary sensitivity study evaluated \(3\times3\), \(4\times4\),
\(5\times5\), and \(6\times6\) SOM architectures across five random
initializations. Increasing the number of neurons reduced quantization error
but simultaneously increased the fraction of cells classified as SOM
interfaces. Partition stability also depended on architecture and
initialization; in particular, the lowest quantization error did not
correspond to the most stable partition.

The spatial analysis of the 40 highest-priority interface seeds showed that
exact cell selection can vary appreciably between random initializations,
while many of the alternative selections remain spatially close. For the
nominal \(4\times4\) architecture, the mean exact Top-40 Jaccard similarity
was approximately 0.507, whereas the one- and two-cell-layer proximity match
fractions were approximately 0.895 and 0.960, respectively.

These results reinforce the distinction between reproducibility and
robustness. A fixed SOM configuration and random seed can be reproduced
exactly, but this does not imply invariance to alternative architectures or
initializations. The sensitivity study therefore supports treating the
\(4\times4\) SOM as the nominal configuration of the present adaptive
experiments rather than as an optimized architecture. The complete adaptive
CFD workflow was not repeated for every sensitivity configuration, so no
conclusion is drawn regarding which SOM architecture would provide the best
final adaptive-mesh accuracy.

\subsection{Future Work}

The results suggest that further development should focus on
\emph{refinement utility} rather than on further tuning of quantization
error or the Hybrid-v1 coefficient against the existing M160 reference.

The first priority is therefore generalization.

The complete experiment should be repeated at additional Reynolds
numbers while preserving the methodology and indicator definitions used
in the present study. This will determine whether the observed behavior
is specific to the \(Re=100\) flow field or persists as the cavity
dynamics change.

A second priority is evaluation on a different CFD problem. A useful
test should introduce different flow structures while remaining simple
enough to permit systematic mesh verification and controlled reference
solutions. This would provide a stronger test of whether SOM-derived
regime information transfers beyond the lid-driven cavity.

Future work should also examine whether refinement utility can be
estimated more directly. Instead of asking whether a cell has a large
gradient, a large quantization error, or lies near a SOM regime
transition, the relevant question is whether refining that cell or
region produces a measurable improvement in a specified quantity of
interest.

Conceptually, this corresponds to learning or estimating a mapping of
the form

\begin{equation}
\left(
\text{local flow state},
\text{numerical indicators},
\text{mesh state}
\right)
\longrightarrow
\text{expected refinement benefit}.
\end{equation}

Such a formulation would require carefully designed training data and
strict separation between model development and final validation.

Additional investigations should include:

\begin{itemize}

    \item repeating the experiment for additional Reynolds numbers;

    \item testing the methodology on at least one different benchmark
    geometry or flow configuration;

    \item separating spatial and temporal discretization effects through
    independent time-step refinement;

    \item evaluating quantities of interest in addition to global
    velocity-field discrepancies;

    \item investigating repeated adaptation cycles rather than only
    one-shot refinement;

    \item extending the SOM architecture and initialization sensitivity
    analysis to complete adaptive CFD calculations, while also examining
    alternative feature sets through controlled ablation studies;

    \item comparing SOM-derived information with established
    error-estimation or adaptation methods;

    \item measuring total adaptive-workflow cost, including pilot
    simulation, feature extraction, machine-learning operations,
    remeshing, and final CFD solution;

    \item repeating timing measurements before drawing conclusions about
    computational efficiency.

\end{itemize}

Importantly, future indicator development should avoid optimizing
parameters against the same fine-grid reference subsequently used for
evaluation. New hypotheses should be defined using pilot-solution
information and then tested against independent numerical references.

The longer-term objective is therefore not to replace conventional CFD
adaptation with a SOM. It is to determine whether learned representations
of the physical flow contain information that can complement established
numerical indicators and improve the allocation of computational
resources in a reproducible and generalizable manner.
